\section{Obstruction atlas}

The observed difficulty of finding one additional point is not a single
phenomenon.  It is the product of several mathematically different
obstructions.

\subsection{Surface-lattice capacity}

For a relatively minimal elliptic surface $\pi:S\to\mathbf P^1$ with section,
Shioda--Tate gives
\[
\rank \MW(S)=\rho(S)-2-\sum_v(m_v-1),
\]
where $m_v$ is the number of irreducible components of the fibre at $v$.
For arithmetic genus $d=\chi(\mathcal O_S)$ one has
\[
h^{1,1}(S)=10d,
\qquad
\rank\MW(S)\le10d-2.
\]
Thus a rational elliptic surface has generic rank at most $8$, a K3 elliptic
surface at most $18$ geometrically, and an arithmetic-genus-three surface at
most $28$.  These are generic-rank bounds; they do not bound a single
specialized fibre over $\Q$.

Every reducible fibre consumes Picard rank through its root lattice.  This is
the \emph{reducible-fibre cost}.  A family designed for easy torsion or
isogeny descent can therefore have insufficient capacity before any search is
run.

\subsection{Specialization and visibility}

A fixed independent generic subgroup specializes injectively away from a
controlled exceptional set, but extra points discovered after specialization
can be dependent.  Given a generic lattice with Gram matrix $G_0$ and
candidate extra points with block matrix
\[
G=\begin{pmatrix}G_0&B\\B^T&C\end{pmatrix},
\]
the new directions are measured by the Schur complement
\[
S=C-B^TG_0^{-1}B.
\]
The relevant search height is the norm in the quotient lattice represented by
$S$, not merely the absolute canonical height of a candidate point.

In a non-isotrivial family, specialized canonical heights grow with the height
of the base parameter.  A curve can therefore have genuinely high rank while
its missing generators lie beyond a feasible coordinate search.  This is a
visibility obstruction, not a rank obstruction.

\subsection{Descent and the local--global gap}

A large $2$- or $3$-Selmer group records locally soluble covering classes.
The excess over visible Mordell--Weil rank may lie in $\Sha$.  Consequently,
\[
\text{large Selmer dimension}
\quad\not\Longrightarrow\quad
\text{many rational points}.
\]
The correct construction problem is to make the relevant covering curves
globally soluble and of manageable height, preferably through shared
low-genus geometry.

\subsection{Auxiliary-base arithmetic}

Extra points on fibres are naturally represented by multisections or covering
curves.  A genus-zero auxiliary curve with a rational point is parametrizable.
A genus-one auxiliary curve is useful when its Jacobian has positive rational
rank.  A higher-genus auxiliary curve can have only finitely many rational
points, although finiteness never implies emptiness.

Thus the useful resource is not only generic rank.  It is
\[
\text{forced rank}
+\text{low auxiliary genus}
+\text{positive-rank parameter space}
+\text{low quotient height}.
\]

\subsection{Certification cost}

A lower-bound certificate does not require complete saturation of the full
Mordell--Weil group, but it does require exact independence of the displayed
points.  Finite-field reductions, interval-certified heights, and a second
computer algebra system should be used incrementally.  Search code and
verification code are deliberately separated so that a verifier depends only
on the equation and point list.

\subsection{What the historical graph does not prove}

Record dates are not independent observations from a stationary process.
Each jump reflects a different family, fibration, sieve, point-search method,
amount of computation, and publication decision.  A fitted ``inflection
point'' is therefore not evidence of a universal rank ceiling.  Failure of a
family, a bounded search, or a low-height enumeration yields only a restricted
obstruction unless it is converted into a theorem.
